YS II.44

svadhyayad ista-devata-samprayoga

YS II.45

samadhi-siddhi îsvara-pranidhana

SK 22.

Prakriti mahat tata ahamkara tasmat ca sodasaka gana
sodasakat tasmad api pancabhyah panca bhutani

From nature evolves the Great Principle;
from this evolves the I Principle;
from this evolves the set of sixteen;
from the five of this set of sixteen evolves the five elements.

SK 23.

buddhi adhyavasaya dharma jnana viraga aisvarya
etad rupam sattvikam tasmat viparyastam tamasam

Buddhi is self-certainty.
Virtue, knowledge, dispassion, and power are its manifestations
when the sattva attribute abounds.
And the reverse of these, when the tamas attribute abounds.

SK 24.

ahamkara abhimana tasmat pravartate dvividha sarga eva
ekadasaka gana ca tanmatra pancaka

Ahamkara is self-assertion;
from that proceeds a two-fold evolution
the set of eleven and the five-fold primary elements.

YS II.46

sthira-sukham asanam

YS II.47

prayatna-saithilya-ananta-samapattibhyam

YS II.48

tato dvandva-anabhighata

SK 25.

vaikrita ahamkara ekadasaka sattvika pravartate
tanmatra bhutade sa tamasa taijasa ubhayam

The set of eleven sattvic indriyas proceeds from the Vaikriti form of I-Principle;
the set of five tamasic tanmatras proceed from the Bhutadi form of I-Principle.
From the Taijasa form of I-Principle proceed both of them.

SK 26.

buddhi indriyani akhyani caksu srotram ghrana rasana tvak
karma indriyani ahu vak pani pada payu upastha

Organs of knowledge are called the Eye, the Ear, the Nose, the Tongue, the Skin.
Organs of action are called the Speech, the Hand, the Feet, the excretory organ and the organ of generation.

YS II.49

tasmin sati svasa-prasvasayo gati-viccheda pranayama

YS II.50

bahya-abhyantara-stambha-vritti desa-kala-sankhyabhi-paridrsta dirgha-suksma

YS II.51

bahya-abhyantara-visaya-aksepi caturtha

SK 27.

atra mana ubhayatmakam sankalpakam ca sadharmaya indriyam
nanatvam bahya bheda ca guna-parinama-visesa

Of these (sense organs), the Mind possesses the nature of both.
It is the deliberating principle, and is also called a sense organ
since it posseses properties common to the sense organs.
Its multifariousness and also its external diversities are owing
to special modifications of the Attributes.

SK 28.

pancanam rupadisu alocanamatram isyate
pancanam vritti vacana adana viharana utsarga ca ananda

The function of the five in respect to form and the rest,
is considered to be mere observation.
Speech, manipulation, locomotion, excretion and gratification
are the functions of the other five.

YS II.52

tata ksiyate prakasa-avaranam

YS II.53

dharanasu ca yogyata manasa

YS II.54

sva-visaya-asamprayoge cittasya svarupa-anukara ive indriyanam pratyahara

YS II.55

tata parama vasyata indriyanam

YS III.1

desa-bandha cittasya dharana

YS III.2

tatra pratyaya-eka-tanata dhyanam

YS III.3

tad evartha-matra-nirbhasa svarupa-sunyam iva samadhi

YS III.4

trayam ekatra samyama

YS III.5

taj-jayat prajna-aloka

YS III.6

tasya bhumisu viniyoga

YS III.7

trayam antar-angam purvebhya

YS III.8

tad api bahir-angam nirbijasya

SK 29.

trayasya svalaksanyam vritti sa esa asamanya bhavati
samanya karana-vritti pranadhyah-vayavah-panca

Of the three internal instruments, their own characteristics are their functions;
this is peculiar to each. The common modification of the instruments is the five airs
such as prana and the rest.

The concept is the free [actuality], as the substantial power that is for itself,
and it is the totality, since each of the moments is the whole that it is,
and each is posited as an undivided unity with it.
So, in its identity with itself, it is what is determinate in and for itself.

The way the concept proceeds is no longer passing over or shining in an other.
It is instead development since what are differentiated are at the same time
immediately posited as identical with one another and with the whole,
each being the determinacy that it is as a free being of the whole concept.

The doctrine of the concept is divided into the doctrine of

(1) the subjective or formal concept,
(2) the concept as determined to immediacy, or the objectivity,
(3) the idea, the subject-object, the unity of the concept and objectivity,
the absolute truth.

A. THE SUBJECTIVE CONCEPT

YS III.9

vyutthana-nirodha-samskarayor abhibhava-pradur-bhavau nirodha-kshana-cittanvayo nirodha-parinama

YS III.10

tasya prasanta-vahita samskarat

Manas, the understanding, first the faculty for the cognition of the general (of rules)

a. Concept as such

The concept as such contains the moments of

universality
(as the free sameness with itself in its determinacy),
particularity
(as the determinacy in which the universal remains the same as itself, unalloyed),
and individuality
(as the reflection-in-itself of the determinacies of universality and particularity,
the negative unity with itself that is the determinate in and for itself
and at the same time identical with itself or universal).

The concept is what is utterly concrete since the negative unity with itself
(as being-determined-in-and-for-itself which is the individuality)
itself makes up its relation to itself, the universality.
To this extent, the moments of the concept cannot be detached from one another;
the determinations of reflection are supposed to be grasped
and to be valid each for itself,
detached from the opposed determination.
Since, however, their identity is posited in the concept,
each of its moments can be immediately grasped only
on the basis of and with the others.

SK 30.

drste catustayasya tu yugapat vrttih tasya kramasasca nirdista
tatha api adrste trayasya vrttih tat purvika

Of all the four, the functions are said to be simultaneous and also successive
with regard to the seen objects; with regard to the unseen objects
the functions of the three are preceeded by that.

The moment of individuality first posits the moments of the concept as differences,
since it is the concept's negative reflection-in-itself.
Thus it is initially the free differentiating of the concept as the first negation,
by means of which the determinacy of the concept is posited, but posited as particularity.
That is to say, first, that the moments differentiated have
the determinacy of conceptual moments only opposite one another and,
second, that their identity (that the one is the other) is equally posited.
This posited particularity of the concept is the judgment.

YS III.11

sarva-arthata-ekagratayo kshayodayau cittasya samadhi-parinama

The logical form of all judgments consists of the objective unity
of the apperception of the concepts contained therein

Ahamkara, the determinative power of judgment,
second the faculty for the subsumption of the particular under the general

b. The judgment

The judgment is the concept in its particularity as the differentiating
relation of its moments, which are posited as being for themselves and,
at the same time, as identical with themselves, not with one another.

Judgment is usually taken in the subjective sense as an operation and form
that surfaces merely in self-conscious thinking.
This difference, however, is not yet on hand in the logical [sphere, where]
judgment is supposed to be taken in the completely universal sense:
all things are a judgment;
they are individuals which are a universality or inner nature in themselves,
or a universal that is individuated. The universality and individuality
distinguish themselves in them [the things] but are at the same time identical.

The standpoint of the judgment is finitude, and from this standpoint the
finitude of things consists in the fact that they are a judgment,
that their existence and their universal nature (their body and their soul) are
certainly unified (otherwise the things would be nothing), but that these,
their moments, are both already diverse and generally able to be separated.

In the abstract judgment 'the individual is the universal',
the subject relates itself negatively to itself and, as such,
is the immediately concrete, while the predicate is, by contrast,
the abstract, indeterminate, the universal.
But since they are joined by 'is', the predicate in its universality
must also contain the determinacy of the subject
and it [that determinacy] is the particularity
and the latter is the posited identity of the subject and predicate.
As thus indifferent to this difference of form, it is the content.

As far as the more precise determinacy of subject and predicate is concerned,
the former, as the negative relation to itself, is the
underlying fixity in which the predicate has its subsistence
and is in an ideal way (it inheres in the subject).
Moreover, since the subject is generally and immediately concrete,
the determinate content of the predicate is only one of the many
determinacies of the subject and the latter is richer and broader than the predicate.
Conversely, the predicate, as the universal subsisting for itself and
indifferent to whether this subject is or not, goes beyond the subject,
subsumes the subject under it, and is, for its part, broader than the subject.
The determinate content of the predicate (see preceding section)
alone makes up the identity of both.

Subject, predicate, and the determinate content or the identity [of them]
are initially posited in the judgment, in their relation, as themselves diverse,
falling outside one another.
But in themselves, in terms of the concept, they are identical,
since the concrete totality of the subject is this,
not to be some sort of indeterminate manifold, but instead individuality alone,
the particular and universal in an identity, and precisely this unity is the predicate.
In the copula, furthermore, the identity of the subject and predicate is
of course posited but initially only as the abstract 'is'.
In keeping with this identity, the subject is also to be posited in
the determination of the predicate, by means of which the latter also acquires
the determination of the subject and the copula is fulfilled.
This is the further determination of the judgment,
by means of the copula full of content, into the syllogism.
But first, in terms of the judgment, there is the further determination of it,
the determining of the initially abstract, sensory universality into
a set of all, genus, and species and into the developed universality of the concept.

SK 31.

svam svam vrittim pratipadyante paraspara akuta hetukam
purushartha eva hetu na kenacit karanam karyate

Instruments enter into their respective functions being incited by mutual impulse.
The meaning of the Spirit is the sole motive (for the activities of the instruments).
By none whatsoever is an instrument made to act.

In this way, subject and predicate are each themselves the entire judgment.
The immediate constitution of the subject shows itself at first as the mediating ground
between the individuality of the actual and its universality, as the ground of the judgment.
What has in fact been posited is the unity of the subject and the predicate, as the concept itself;
it is the fulfilment of the empty 'is', the copula,
and since its moments are at the same time differentiated as subject and predicate,
it is posited as their unity, as the relation mediating them - the syllogism.

YS III.12

tata punashantoditau tulya-pratyayau cittasya-ekagrata-parinama

Buddhi, reason, third the faculty for the determination of the particular
from the general (for the derivation from principles)

c. The syllogism

The syllogism is the unity of the concept and the judgment;
it is the concept as the simple identity
(into which the judgment's differences of form have gone back),
and [it is] judgment insofar as it is posited at the same time in reality,
namely, in the difference of its determinations.
The syllogism is what is rational and everything rational.

The immediate syllogism is such that the determinations of the concept
stand opposite one another in an external connection as abstract determinations,
so that the two extremes [are] the individuality and universality,
but the concept, as the middle joining the two together,
is likewise only the abstract particularity.
The extremes are accordingly posited as subsisting for themselves,
as indifferent to one another as they are to the middle [term that joins them].
This syllogism is thus rational but non-conceptual;
it is the formal syllogism of the understanding.
In it the subject is joined together with another determinacy;
or through this mediation the universal subsumes a subject external to it.
In a rational syllogism, by contrast, the subject joins itself
together with itself by means of this mediation.
It is only a subject in this way, or the subject is only in itself
the syllogism of reason.

The syllogism has been taken in terms of the differences contained in it
and the universal result of the course of those differences is that
these differences and the concept's manner of being-outside-itself
sublate themselves in it. Indeed,
(1) each of the moments themselves has demonstrated itself to be
the totality of the moments and, hence, to be the entire syllogism;
thus, they are in themselves identical.
(2) The negation of their differences and their mediation constitutes the
manner of being-for-itself, such that it is one and the same universal
that is in these forms and is accordingly also posited as their identity.
In this ideality of the moments, inferring acquires the determination of
essentially containing the negation of the determinacies
by means of which it runs its course and, with this,
the determination of being a mediation by way of sublating the mediation
and a manner of joining the subject together, not with another,
but with the sublated other, with itself.

SK 32.

karanam trayodasavidham tad aharana dharana prakasakaram
tasya karyam ca dasadha aharyam dharyam prakasyam ca

Instruments are of thirteen kinds performing the functions of
seizing, sustaining and illuminating.
Its objects are of ten kinds, vis
the seized, the sustained, and the illumined.

This realization of the concept,in which
the universal is this singular totality that has returned into itself
and whose differences are equally this totality,
which has determined itself to be an immediate unity by sublating mediation,
this realization of the concept is - the object.

B. THE OBJECT

The object is immediate being by virtue of
the indifference towards the difference
that has sublated itself in it.
It is in itself the totality and, at the same time,
since this identity is the identity of the moments
but an identity that only is in itself,
it is just as indifferent to its immediate unity.
It breaks down into differentiated [moments],
each of which is itself the totality.
The object is thus the absolute contradiction of
the complete self-sufficiency of the manifold and
the equally complete lack of self-sufficiency of
the differentiated [moments].

YS III.13

etena bhutendriyesu dharma-laksana-vastha-parinama vyakhyata

a. Mechanism

The object, taken first in its immediacy, is
(1) the concept only in itself,
it has the concept at first as something subjective outside it, and
every determinacy is posited as an external determinacy.
As the unity of differences, it is thus something composite, an aggregate,
and the effect on another remains an external relation: formal mechanism.
In this relation and lack of self-sufficiency, the objects remain equally
self-sufficient, resistant, external to one another.

Only insofar as the object is self-sufficient (see the preceding section)
does it have the lack of self-sufficiency in terms of which it suffers violence.
Insofar as the object is the posited concept in itself,
neither of these determinations sublates itself
in its other determination;
instead the object joins itself together with itself
through the negation of itself, through its lack of self-sufficiency,
and only then is it self-sufficient.
Thus, at the same time, in the difference from externality,
and in its self-sufficiency negating this externality,
it [the object] is the negative unity with itself, centrality,
subjectivity in which it is itself directed and related to the external.
The latter is equally centred in itself and, in that, just as much related
to the other centre, having its centrality just as much in the other.
[Hence, the object in the second place is]
(2) a differentiated mechanism (fall, desire, social drive, and the like).
The development of this relationship forms the syllogistic inference
that the immanent negativity as the central individuality of an object
(the abstract centre) relates itself to objects lacking self-sufficiency as
the other extreme, relating to them through a middle [term] that
unifies the objects' centrality and lack of self-sufficiency,
the relative centre. [Hence, the object is]
(3) absolute mechanism.

The syllogism that has been given here (I - P - U) is
a triad of syllogistic inferences.
The flawed individuality of the objects lacking self-sufficiency,
in which the formal mechanism is at home, is, in keeping with
its lack of self-sufficiency, just as much the external universality.
These objects are thus the middle also between the absolute and
the relative centre (the form of the syllogism: U - I - P).
For it is by means of this lack of self-sufficiency
that those two are separated and are extremes
just as they are related to one another.
So, too, the absolute centrality as the substantial universal
(the gravity that remains identical),
which as the pure negativity also encapsulates in itself the individuality,
is the mediating factor between the relative centre
and the objects lacking self-sufficiency;
([thus amounting toJ the form of the inference P - U - I)
and, to be sure, just as essential in terms of the immanent individuality
where it functions to separate, as it is in terms of the universality
as the identical cohesion and as the undisturbed being-in~itself.

SK 33.

antahkaranam trividham bahyam dasadha trayasya visayakhyam
bahyam sampratakalam abhyantaram karanam trikalam

The internal instruments are three-fold.
The external are ten-fold; they are called the objects of the three (internal instruments).
The external instruments function at the present time and
the internal instruments function at all the three times.

The immediacy of concrete existence that objects have in absolute mechanism
is in itself negated by the fact that their self-sufficiency is mediated by
their relations to one another, hence, through their lack of self-sufficiency.
Thus, the object must be posited as differentiated in its concrete existence,
opposite its other.

YS III.14

shantoditavyapadesya-dharmanupati dharmi

b. Chemism

The differentiated object has an immanent determinacy constituting
its nature and in that determinacy it has concrete existence.
But as the posited totality of the concept, it is the contradiction of
this its totality and the determinacy of its concrete existence;
it is thus the [process of] striving to sublate this contradiction
and make its existence equal to the concept.

The chemical process thus has as its product the neutral dimension of
these strung-out extremes, a neutral dimension which these extremes are
in themselves; by means of the differentiation of the objects (the particularization),
the concept, the concrete universal, joins itself
with the individuality, the product, and so merely
with itself equally contained in this process are the other syllogisms;
the individuality, as activity, is likewise the mediating factor just like the concrete universal,
the essence of the strung-out extremes, which enters into existence in the product.

Chemism, as the reflexive relationship of objectivity, still presupposes,
together with the differentiated nature of the objects,
the immediate self-sufficiency of those same objects.
The process is that of passing back and forth from one form into the other,
forms determinate properties that the extremes had opposite one another are sublated.
This is, indeed, in keeping with the concept;
but the animating principle of differentiating does not exist
concretely in it since it has sunk back into immediacy.
For this reason, the neutral dimension is a separable dimension.
Yet the judging principle that severs the neutral dimension into
differentiated extremes and gives the undifferentiated objects in general
their difference and animation opposite an other falls outside that first process,
and so does the process as the separation that strings things out.

SK 34.

Tesam panca buddhi visesa-avisesa-visaya
Vak sabda-visaya-bhavati sesani tu pancavisayani

Of these, the five organs of knowledge have as their objects
both the specific and the general.
Speech has speech as its object;
the rest have all the five as its object.

The externality of these two processes,
the reduction of what are differentiated to something neutral and
the differentiation of the undifferentiated or neutral,
which allows them to appear as self-sufficient opposite one another,
shows its finitude in passing over into products in which they are sublated.
Conversely, the process presents the presupposed immediacy
of the differentiated objects as a vacuous immediacy.
By means of this negation of externality and immediacy,
into which the concept as object was immersed,
it is posited freely and for itself opposite
that externality and immediacy - as purpose.

YS III.15

kramanyatvam parinamanyatve hetu

c. Teleology

Purpose is the concept that is for itself and that has entered into
a free concrete existence via the negation of immediate objectivity.
It is determined as something subjective, in that this negation
initially is abstract and thus objectivity also only stands
over against it [i.e. the purpose] at first.
In contrast to the totality of the concept, however,
this determinacy of the subjectivity is one-sided and, indeed,
for it [the purpose] itself, since all determinacy has posited
itself as sublated in it.
Thus, too, for it [the purpose] the presupposed object is
only an ideal, in itself vacuous reality.
As this contradiction of its identity with itself opposite the negation and the
opposition posited in it, it is itself the sublating, the activity of so negating
the opposition that it posits it as identical with itself.
This is the process of realizing the purpose in which,
by rendering itself something other than its subjectivity and
objectifying itself, it has sublated the difference of both,
has joined itself together only with itself and has preserved itself.

The teleological relation in its immediacy is initially the external purposiveness,
and the concept is opposite the object which is something presupposed.
The purpose is thus finite, partly in terms of the content,
partly in terms of the fact that it has an external condition
in an extant object as the material of its realization.
To this extent, its self-determination is merely formal.
The immediacy entails, more precisely, that the particularity
(as a determination of form, the subjectivity of the purpose)
appears as reflected in itself, the content as distinct from
the totality of the form, the subjectivity in itself, the concept.
This diversity constitutes the finitude of the purpose within itself.
The content is, by this means, as limited, contingent, and
given as the object is something particular and extant.

The teleological relation is the syllogism in which the subjective purpose
joins itself together with the objectivity external to it
through a middle term that is the unity of the two,
both as the purposive activity and as the objectivity
immediately posited under the purpose, the means.

1. The subjective purpose is the syllogism in which the universal concept
joins together with individuality by means of particularity;
such that this [individuality] as the self~determination judges.
That is to say, this individuality both particularizes that still indeterminate universal,
making it a determinate content, and also posits the opposition of subjectivity and objectivity.
It [this individuality] is, in itself, at the same time the return
into itself since it determines the concept's subjectivity
(presupposed as something opposite the objectivity)
to be deficient in relation to the totality that has joined together
with itself and since at the same time it thereby turns outward.

2. This activity turned outward is the individuality that,
in the subjective purpose, is identical to the particularity
in which, next to the content, the external objectivity is also included.
As such, this activity relates at the outset immediately to the object and
takes control of it as a means.
The concept is this immediate power because it is
the negativity identical with itself,
the being of the object is thoroughly determined only as something ideal.
The entire middle term is now this inner power of the concept as activity,
with which the object is immediately unified as means and under which it stands.

3. The purposive activity with its means is still directed outward,
since the purpose is also not identical with the object;
thus it must first be mediated with the object.
The means, as the object in this second premise, is in
immediate relation with the other extreme of the syllogism,
the objectivity as presupposed, the material.
This relation is the sphere of the mechanism and chemism
now serving the purpose that is their truth and free concept.
That the subjective purpose, as the power of these processes in which the
objective dimension rubs up against itself and sublates itself,
keeps itself outside them and is what preserves itself in them;
this is the cunning of reason.

Thus, the realized purpose is the posited unity of
the subjective and the objective dimensions.
This unity, however, is essentially determined in such a way that
the subjective and objective dimensions are neutralized
and sublated only with respect to their one-sidedness,
while the objective dimension is subjected and made to conform
to the purpose as the free concept and, thereby, to the power over it.
The purpose preserves itself against and in the objective dimension
because, in addition to being the one-sided subjective dimension (the particular),
it is also the concrete universal, the identity of both, that is in itself.
This universal, that as simple is reflected in itself, is the content that
remains the same through all three termini [terms] of the syllogism and their movement.

In the finite purposiveness, however, the purpose carried out is also
something as internally broken as was the middle term and the initial purpose.
What has come about is thus only a form posited externally in the material
found before it, a form that, on account of the restricted content of
the purpose, is likewise a contingent determination.
The purpose attained is thus only an object that is also
in turn a means or material for other purposes and so on ad infinitum.

SK 35.

yasmat buddhi santahkarana sarvam visayam avagahate
tasmat trividham karanam dvari sesani dvarani

Since buddhi along with the other internal instruments
comprehends all of the objects these three instruments are
like the warders while the rest are like the gates.

What happens, however, in the process of realizing
the purpose in itself is that the one-sided subjectivity and
the semblance of objective self-sufficiency on hand opposite it are sublated.
In seizing the means, the concept posits itself as the object's essence as
it is in itself, in the mechanical and chemical process,
the self-sufficiency of the object has already evaporated in itself and
in the course it takes under the dominance of the purpose,
the semblance of that self-sufficiency,
the negative dimension opposite the concept, sublates itself.

Yet this object is immediately already posited as vacuous in itself,
as only ideal by virtue of the fact that the executed purpose is
determined only as means and material.
With this, the opposition of content and form has vanished as well.
Since the purpose, by sublating the formal determinations,
joins itself together with itself,
the form is posited as identical with itself, thus as content,
so that the concept as the activity of the form has only itself as content.
It is thus posited through this process generally what the concept of the purpose was:
the unity, being in itself, of the subjective and the objective dimensions
now posited as being for itself - the idea.

C. THE IDEA

The idea is the true in and for itself,
the absolute unity of the concept and objectivity.
Its ideal content is none other than the concept in its determinations.
Its real content is only its exhibition, an exhibition
that it provides for itself in the form of external existence and,
with this shape incorporated into the concept's ideality and in its power,
the concept thus preserves itself in that exhibition.

The idea can be grasped as reason
(this is the genuine philosophical meaning of reason),
further as subject-object, as the unity of the ideal and the real,
of the finite and the infinite, of the soul and the body,
as the possibility that has its actuality in itself,
as that the nature of which can only be conceived as existing,
and so forth, because in it [the idea] all relationships of
the understanding are contained,
but in their infinite return and identity in themselves.

The idea is essentially a process since its identity is
that of the absolute and free concept only insofar as
it is the absolute negativity and thus dialectical.
It is the course in which the concept as the universality that is individuality
determines itself to be objectivity and to be the opposite of objectivity,
and in which this externality that has the concept as its substance leads
itself back into subjectivity through its immanent dialectic.

YS III.16

parinama-traya-samyamad atitanagata-jnanam

a. Life

The immediate idea is life.
The concept is realized as the soul in a body;
the soul is the immediate, self-referring universality of the body's externality,
just as much as it is the body's particularization, so that the body expresses
no other differences than the determination of the concept, and finally
it is the individuality as infinite negativity;
the dialectic of the body's objectivity,
[the factors of which are] outside one another,
an objectivity that is led back into subjectivity from the semblance of self-sufficient
subsistence, so that all members are reciprocally momentary means as
much as momentary purposes, while life, inasmuch as it is the inceptive
particularization, results in itself as the negative unity that is for itself and,
in the dialectic of embodiment, joins itself together only with itself.

Life is thus essentially a living entity and, with regard to its immediacy,
this individual living entity.
In this sphere, finitude has the determination that soul and body are separable,
on account of the immediacy of the idea;
this constitutes the mortality of the living.
But those two sides of the idea are diverse component parts
only insofar as it is dead.

The living is the syllogism, whose moments are systems and syllogisms
in themselves which, however, are active syllogisms, processes, and
in the subjective unity of the living, they are only one process.
The living is thus the process of its coming to closure together with itself,
that runs its course by means of three processes.

1. The first is the process of the living within itself,
in which it divides itself in itself and makes its
corporal condition its object, its inorganic nature.
For its part, this inorganic side, as the relatively external,
enters into the difference and opposition of its moments that reciprocally
surrender themselves, the one assimilating the other to itself, and
preserve themselves in the process of producing themselves.
This activity of the members, however, is only one activity of
the subject, the activity into which its productions go back,
so that through that activity only the subject is produced;
it merely reproduces itself.

2. But the judgment of the concept proceeds freely to release from itself
the objective dimension as a self-sufficient totality.
The negative relation of the living to itself, as immediate individuality,
presupposes an inorganic nature standing over against it.
Since this negative aspect of itself is just as much a moment of
the concept of the living itself, it is thus in the latter
(the at once concrete universality) as a lack.
The dialectic, through which the object as something in itself vacuous sublates itself,
is the activity of the living entity certain of itself that accordingly
preserves, develops, and objectifies itself in this process opposite an inorganic nature.

3. In the initial stage of its process, the living individual
behaves as a subject and concept in itself.
Through its second stage, it assimilates its external objectivity to itself and
thus posits in itself the real determinacy.
As a result, it is now in itself the genus, substantial universality.
The particularization of the latter is the relation of
the subject to another subject of its genus and
the judgment is the relationship of the genus to these
determinate individuals standing opposite one another:
the difference of the sexes.

The process of the genus brings this [genus] to
the point of being-for-itself.
Because life is still the immediate idea,
the product of the process breaks down into two sides.
On the one side, the living individual in general,
at first presupposed as immediate,
emerges now as something mediated and produced.
On the other side, however, the living individuality that,
on account of its initial immediacy,
behaves negatively towards the universality,
perishes in this [universality] as the power.

SK 36.

ete manas ahamkara pradipa-kalpa paraspara vilaksana
guna-visesa krstnam prakasya purushasya-artham buddhau prayacchati

These (external instruments along with manas and ahamkara)
which are characteristic-wise different from one another
and are different modifications of the attributes
and which resemble a lamp (in action) illuminating all (their respective objects)
present them to the Buddhi for the purpose of the Spirit,

By this means, however, the idea of life has not only freed itself from just
any (particular) immediate 'this', but from this initial immediacy altogether.
In this way, it comes to itself, to its truth,
entering into concrete existence as the free genus for itself.
The death of the merely immediate, individual living thing is the Spirit emerging.

The idea concretely exists freely for itself insofar as universality is the element
in which it exists concretely or insofar as it is objectivity itself as the concept;
[that is to say,] the idea has itself for an object.
Its subjectivity, determined as universality, is pure differentiating within it
intuiting that keeps itself in this identical universality.
But, as a differentiating in a determinate way, it is the further judgment of
thrusting itself as a totality away from itself and, indeed,
initially presupposing itself as the external universe.
These are two judgments that are in themselves identical but
not yet posited as identical.

The relation of these two ideas that are identical in themselves or as life is thus
the relative relation that makes up the determination of finitude in this sphere.
It is the relationship of reflection, since the differentiation of
the idea in it [the idea] itself is only the first judgment,
the presupposing is not yet a positing, and thus,
for the subjective idea, the objective dimension is the extant immediate world or
the idea as life in the appearance of individual concrete existence.
At the same time, insofar as this judgment is a pure differentiating
within it [the idea] itself (see the preceding section),
the idea is for itself both itself and its other.
Thus it is the certainty of being in itself
the identity of this objective world with it.
Reason comes to the world with the absolute faith
in its capacity to posit the identity and
elevate its certainty to truth, and
with the drive to posit as also vacuous for it
that opposition that is in itself vacuous.

In general, this process is knowing.
In it, in one activity, the opposition,
the one-sidedness of subjectivity together with
the one-sidedness of objectivity,
is sublated in itself.
But this process of sublating takes place at the outset only in itself.
The process as such is thus itself immediately beset with the finitude
of this sphere and falls apart into the twofold, diversely posited
movement of the drive.
[In one respect,] it is the drive to sublate
the one-sidedness of the subjectivity of the idea
by taking up into itself the world that is,
taking it up into subjective representing and thinking,
and to fill out the abstract certainty of itself
with this objectivity as content, an objectivity that thus counts as true.
Conversely, it is the drive to sublate
the one-sidedness of the objective world
that here accordingly, by contrast, counts as a semblance,
a collection of contingencies and shapes vacuous in themselves,
and to determine and mould it through the inner dimension of the subjective,
that counts here as the objective, as what truly is.
The former is the drive of knowledge to truth, knowing as such,
the theoretical (activity);
the latter is the drive of the good to bring itself about, willing,
the practical activity of the idea.

YS III.17

sabda-artha-pratyayanam itaretara-adhyasat sankaras tat-pravibhaga-samyamat sarva-bhuta-ruta-jnanam

YS III.18

samskara-saksat-karanat purva-jati-jnanam

YS III.19

pratyayasya para-citta-jnanam

YS III.20

na ca tat salambanam tasya-avisayi-bhutatvat

b. Knowing

The universal finitude of knowing that lies in the first judgment,
the presupposition of the opposition, which its very action contradicts,
specifies itself more precisely in its own idea in this direction,
that its moments receive the form of diversity from one another and,
since those moments are in fact complete, they come to stand in
the relationship of reflection, not of the concept, to one another.
The assimilation of the material as something given thus appears as
a way of taking it up into conceptual determinations that at the same time
remain external to it, determinations that likewise display themselves
opposite one another as diverse.
It is reason active as understanding.
The truth that this knowing comes to is thus likewise only finite;
the infinite truth of the concept is fixed as a goal that is only in itself,
something beyond this knowing.
But in its external action, it stands under the guidance of the concept,
and conceptual determinations make up the inner thread of the progression.

Because it presupposes what is differentiated as a being that is
found to be already on hand, standing opposite it
(the manifold facts of external nature or of consciousness),
finite knowing has (1) the formal identity or
the abstraction of universality as the form of its activity at the outset.
This activity thus consists in dissolving the given concrete dimension,
individuating its differences, and giving them the form of abstract universality;
or in leaving the concrete dimension as the ground and,
through abstraction from the particularities that seem inessential,
extracting a concrete universal, the genus or the force and the law.
Such is the analytic method.

This universality is (2) also a determinate one.
The activity here proceeds according to the moments of
the concept that, in finite knowing, is not in its infinity
but is the understandable , determinate concept instead.
Taking up the object into the forms of
the latter concept is the synthetic method.

(AA) Knowing initially puts the object into the form of
the determinate concept in general so that, by this means,
its genus and its universal determinacy are posited.
The respective object is the definition.

Its material and justification are procured by the analytic method.
The determinacy is, nevertheless, supposed to be only a characteristic,
that is to say, something to assist merely subjective knowing
that is external to the object.

(BB) The account of the second moment of the concept,
the determinacy of the universal as particularization,
is given by the division in terms of some sort of external aspect.

(CC) In the concrete individuality (such that the simple determinacy in
the definition is construed as a relationship),
the object is a synthetic relation of differentiated determinations - a theorem.

Because they are diverse, their identity is a mediated identity.
The process of supplying the material that constitutes
the middle members is the construction;
and the mediation itself, out of which the necessity of
that relation for knowing goes forth, is the proof.

SK 37.

yasmat buddhi purusasya upabho-gam sarvam pratisadhyati
sa eva ca suksma pradhana-purusantaram visinasti

Because it is the Buddhi that accomplishes the experiences
with regard to all objects to the Purusha.
It is that again that discriminates the subtle difference
between the Pradhana and the Purusha.

The necessity which finite knowing produces in a proof
is initially an external necessity,
determined only for the subjective discernment.
But in the necessity as such, it has itself left behind its presupposition
and point of departure, the finding and givenness of its content.
The necessity as such is, in itself, the concept relating itself to itself.
The subjective idea has thus, in itself, come to what is determined in and for itself,
what is not given, and is thus immanent to it as the subject.
As such, it passes over into the idea of willing.

The subjective idea - as what is determinate in and for itself,
the simple, self-same content - is the good.
Its drive of realizing itself inverts the relationship
that holds relative to the idea of the true, and is bent on determining,
in terms of its purpose, the world that it finds.

This willing is, on the one hand, certain of the vacuousness of the presupposed object
but, on the other hand, as finite, it at the same time presupposes
both the purpose of the good as a merely subjective idea
and the independence of the object.

The finitude of this activity is thus the contradiction that, in the self-
contradicting determinations of the objective world, the purpose of the good
is both carried out and not carried out, and that it is posited as something
inessential just as much as something essential, as something actual and at
the same time as merely possible. This contradiction presents itself as the
endless progression in the actualization of the good, that is therein established
merely as an ought. Formally, however, this contradiction disappears in
that the activity sublates the subjectivity of the purpose and thereby the
objectivity, the opposition through which both are finite, and not only the
one-sidedness of this subjectivity but subjectivity in general; another such
subjectivity, that is to say, a new generation of the opposi tion, is not distinct
from what was supposed to be an earlier one. This return into itself is at
the same time the recollection of the content into itself, which
is the good and the identity in itself  of both
sides, - the recollection of the presupposition of the theoretical stance,
that the object is what is substantial in itself and true.

The truth of the good is, by this means, posited as
the unity of the theoretical and practical idea,
[the notion] that the good has been attained in and for itself,
that the objective world is thus in and for itself the idea precisely as
it [the idea] at the same time eternally posits itself as purpose and
through activity produces its actuality.
This life, having come back to itself from
the differentiation and finitude of knowing,
and having become identical with the concept
through the activity of the concept,
is the speculative or absolute idea.

YS III.21

kaya-rupa-samyamat tad-grahya-sakti-stambhe caksu-prakasa-asamprayoge ‘ntardhanam

YS III.22

etena sabda-adi-antardhanam uktam

c. Absolute Idea

The idea as the unity of the subjective and the objective idea is
the concept of the idea, for which the idea as such is the object,
for which it is the object - an object into which all determinations have gone together.
This unity is accordingly the absolute and entire truth,
the idea thinking itself, and here, indeed, as thinking, as the logical idea.

The absolute idea is for itself, since in it there is no transition or presupposing
and no determinacy at all that is not fluid and transparent;
it is the pure form of the concept that intuits its content as itself.
It is content for itself insofar as it is the ideal differentiating of
itself from itself, and one side of what has been differentiated is
the identity with itself, in which, however, the totality of the form is
contained as the system of the determinations of the content.
This content is the system of the logical.
Nothing remains here of the idea, as form, but the method of this content -
the determinate knowledge of the validity of its moments.

The moments of the speculative method are

(A) The beginning, which is being or the immediate;
for itself for the simple reason that it is the beginning.
From the vantage point of the speculative idea, however,
it is the speculative idea's self-determining which,
as the absolute negativity or movement of the concept,
judges and posits itself as the negative of itself.

Being, which from the vantage point of the beginning as such
appears as abstract affirmation, is thus instead the negation,
positedness, being-mediated in general and being pre-supposed.
But as the negation of the concept that is simply identical
with itself in its otherness and is the certainty of itself,
it is the concept not yet posited as concept or, in other words,
it is the concept in itself.
For that reason, as the still undetermined concept,
i.e. the concept determined only in itself or immediately,
this being is just as much the universal.

(B) The progression is the posited judgment of the idea.
The immediate universal, as the concept in itself, is
the dialectic of reducing, within itself,
its immediacy and universality to a moment.
It is accordingly the negative [aspect] of the beginning or
the first [moment] posited in its determinacy;
it is for something, the relation of what has been differentiated
- the moment of reflection.

The abstract form of the progression within the stage of being is
[to be] an other and a passing over into an other;
in the stage of the essence, it is the shining in something opposite;
in the stage of the concept, it is the differentiated status
of the individual from the universality which continues itself as such
in what is differentiated from it and is as an identity with the latter.

In the second sphere, the concept at first being in itself came to
shine forth and is thus in itself already the idea.
The development of this sphere becomes the return to the first,
just as the development of the first sphere is a transition into the second.
Only by means of this double movement is justice done to the difference,
since each of the two differentiated factors, each considered in itself,
completes itself so as to form the totality and, in that totality,
puts itself into unity with the other.
Only the fact that both sublate the one-sidedness
in themselves prevents the unity from becoming one-sided.

The second sphere develops the relation of what has been differentiated
into what the relation is at first, namely a contradiction in the relation itself
in the infinite progression.

(C) This contradiction resolves itself into the end,
where the differentiated is posited as what it is in the concept.
It is the negative of the first, and, as the identity with the latter,
the negativity of itself.
Hence, it is the unity in which these first two, as ideal and as moments,
are sublated, i.e. preserved at the same time.
The concept, starting out from its being-in-itself, thus comes to a close
with itself by means of its difference and the process of sublating that difference.
This concept is the realized concept,
the concept that contains the positedness
of its determinations in its being-for-itself;
it is the idea for which, as the absolutely first (in the method),
this end is at the same time nothing more than the process by which
the semblance that the beginning is something immediate and
it [the idea] a result vanishes;
in other words, this end is the knowledge that the idea is the one totality.

In this way, the method is not an external form but
the soul and concept of the content,
from which it is distinguished only insofar as
the moments of the concept, even in themselves,
in their [respective] determinacy,
come to appear as the totality of the concept.
Insofar as this determinacy or the content, with the form,
leads itself back to the idea, this idea exhibits itself
as the systematic totality which is only one idea,
the particular moments of which are in themselves this same idea
to the same extent that they bring forth the simple being-for-itself
of the idea through the dialectic of the concept.
The science concludes in this way by grasping the concept of
itself as the pure idea, for which the idea is.

The idea, which is for itself, considered in terms of this, its unity with itself,
is the process of intuiting and the idea insofar as it intuits is nature.
As intuiting, however, the idea is posited by external reflection
in a one-sided determination of immediacy or negation.
Yet the absolute freedom of the idea is that it does not merely pass over into life
or let life shine in itself as finite knowing, but instead,
in the absolute truth of itself, resolves to release freely from itself
the moment of its particularity or the first determining and otherness,
the immediate idea, as its reflection, itself as nature.
